\section{Organic and Mechanical Thinking}

\label{sec:OrganicandMechanicalThinking}

Organic and mechanical thinking and their impact on Tradition in the West.

\begin{quotex}
\begin{verse}
To see a world in a grain of sand,\\
And a heaven in a wild flower,\\
Hold infinity in the palm of your hand,\\
And eternity in an hour.
\end{verse}
\flright{\textsc{William Blake}}

\end{quotex}
\paragraph{Intellectual Intuition}
In his \textit{Lectures on Divine Humanity}, \textbf{Vladimir Solovyov} distinguishes between organic and mechanical thinking, the former related to intellectual intuition and the latter to abstract reasoning. We will relate this to the understanding of caste and in a later post show how this relates to the Medieval project of implementing the tri-functional society.

Rational thought abstracts from the data of sense experience by a process of negation, that is, by excluding specific features of the object of experience. As such, therefore, these abstractions have no positive content and are mental constructs. That is why Kant was able to demonstrate that the process of pure reason can never lead to certain knowledge about the entities of sense experience.

Ideas, on the contrary, are not the results of abstraction, but are inherent to the entities themselves. This means they are neither rationalist abstractions nor sensuous realities, but have an independent reality. Consequently, only a special form of thought activity can know them; this is called intellectual intuition. Abstract thinking is an intermediate stage:

\begin{quotex}
it appears when the mind is strong enough to liberate itself from the exclusive domination of sense perception … but not yet capable of grasping an idea in the fullness and integrity of is actual objective being, of uniting with it inwardly and essentially 

\end{quotex}
Solovyov points to artistic creation as the proof of such intuition, since it is neither the reproduction of sensual phenomena nor general concepts abstracted from that reality. Now Solovyov gets to the heart of the matter:

\begin{quotex}
The truly artistic form or type necessarily requires an inner union of perfect individuality with complete generality or universality. A union of this kind constitutes the essential feature or proper determination of an intellectually intuited idea. 

\end{quotex}
We see this understanding clearly in the poetic mind, as Blake ``sees" the universal in the particular. A similar process is at work in revelation, which itself is often poetical. In revelation, the particular is simultaneously the universal. That is why it is expressed in terms of stories, legends, and myths. Neither the sensuous mind nor the rational mind can see anything beyond the particular events.

\paragraph{Organic Thinking}
Organic thinking is related to intellectual intuition. As such, it is a ``developing" type of thinking, although we prefer to call it thinking in depth or three-dimensional thinking. In other words, it is not ``developing" in the sense that it overturns earlier moments in thought, but rather that it understands ideas in more depth, it sees ever more interconnections of ideas. It is a free activity of mind, and cannot be determined nor ``evolve" strictly as a function of time. Solovyov lists these characteristics:

\begin{itemize}
\item Considers an object in its all-sided wholeness and its inner bond with all other objects 
\item Deduces form within each concept all the others or develops a single concept not the fullness of the whole truth 
\item Perceives or grasps the integral idea of an object 
\end{itemize}
\paragraph{Mechanical Thinking}
Mechanical thinking cannot penetrate into the essence of ideas, so it can only see them in relationship to each other. This type of thinking, therefore, is polemical; it needs to takes sides in a debate, to prefer one idea over another. In the fields of religion or metaphysics, this type of thinking cannot see beyond the dogmas that codify certain insights. It is characterized by:

\begin{itemize}
\item Takes concepts in their abstract separateness 
\item Considers objects under some particular, one-sided definition 
\item Contrasts them one with another in an external manner 
\item Breaks up or differentiates or analyzes immediate reality 
\item Cannot give it a new higher unity 
\item Unable to apprehend the inner truth or meaning (logos) of objects 
\end{itemize}
\paragraph{Agents of Thought}
Solovyov makes it more interesting by identifying those groups with the various groups or castes. In particular, there are ``those who know", the Brahmins or \emph{oratores}, who engage in organic thinking. Solovyov writes:

\begin{quotex}
If this intuition is united with a clear consciousness and is accompanied by reflection, which gives logical definitions of the intuited truth, we have that speculative thinking that characterizes the philosophical art [i.e., metaphysical realization] 

\end{quotex}
Interestingly, Solovyov relates organic thinking to the masses\footnote{\url{https://www.gornahoor.net/?p=5227}}, those having a ``body with a human consciousness". That is, their intellectual intuition is at best virtual, and is thus dependent on their spiritual leaders for its actualization. That relationship explains why the spiritual leaders are so concerned about the care of the poor, the widows, orphans, and so on. In Solovyov's words:

\begin{itemize}
\item \begin{quotex}
If speculative thinking remains in its immediacy and does not clothe its concrete patterns in logical forms, it is the kind of living thought characteristic of people who have not emerged from the unreflective life in their common tribal or national unity 
\end{quotex}
\item \begin{quotex}
Such thinking expresses what is called the folk spirit: 

\end{quotex}
\begin{itemize}
\item \begin{quotex}
Folk creation in art and religion 
\end{quotex}
\item \begin{quotex}
Living development of language 
\end{quotex}
\item \begin{quotex}
Myths superstitions, folk tales, songs, etc 
\end{quotex}
\end{itemize}
\end{itemize}
So we see that the mass is focused on sensory and psychic imagery; the spiritual authority then gives them their universal significance.

Opposed to them are, as Solovyov writes:

\begin{quotex}
Between these two groups are the majority of co-called educated or enlightened people: detached from the worldview of the people, but have not attained an integral philosophical consciousness. These so-called enlightened people are limited to the abstract mechanical thinking that breaks up or differentiates or analyzes immediate reality (and this constitutes the significance and merit of such thinking), but they are not in a position to give it a new, higher unity and union; that is its limitation.

Certainly, it is possible (and in reality, it often happens) that persons of this group, guided in practical life by the ideas of other people's organic thinking in the form of religious beliefs, take the standpoint of the abstract and mechanical intellect in their own theoretical activity. 

\end{quotex}
\paragraph{Tensions}
As long as the masses and the educated people (e.g., the administrators and producers) acknowledge their limitations and follow the spiritual authorities, things function in a Traditional manner. However, there are some inner tensions that threaten this harmony and concord. If the masses reject the authority, there is revolution, and they revert back to their natural materialism and atheism.

For the others, there is obviously a dualism and contradiction in their worldview. Their own thinking is mechanical, which is absolutely necessary and appropriate in their activities in administration and production. This contradiction can only be resolved in an external manner through submission to the deeper understanding of the spiritual authority. In particular, in the history of the West, Christianity was understood in three different ways: organic, mechanical, and superstitious. As the authority of the spiritual leaders was challenged at different points, these disparate viewpoints eventually unraveled in the process of the degeneration of castes, as we have frequently written about.



\flrightit{Posted on 2012-11-29 by Cologero }

\begin{center}* * *\end{center}

\begin{footnotesize}\begin{sffamily}



\texttt{Avery on 2012-11-29 at 09:54 said: }

Aren't Guenon and Evola essentially Platonian, that is, they see objects as attempts to achieve forms that are essentially separate, and discover their meanings in their separateness? And Evola sees the mechanical/artificial as rising above formless, animal nature. Solovyov's philosophy sounds more like Buddhism to me.


\hfill

\texttt{Cologero on 2012-11-29 at 13:17 said: }

Granted, Avery, I've pulled ideas out of their original context, so there is a strong possibility of misunderstanding. Again, specificity is the key to a discussion. Which statements sound ``Buddhist" to you? Where do Guenon and Evola assert what you say about objects (in their own words).


\hfill

\texttt{Avery on 2012-11-30 at 11:48 said: }

Guenon considers Platonic ideals and Aristotelian forms to be ```archetypes' or the essential principles of things" (Reign of Quantity), and comparable to the qualitative Pythagorean numbers. This seems to demonstrate a commitment to separate principles which manifest as different objects.

Evola similarly writes, ``An equality may exist on the plane of a mere social aggregate or of a primordial, almost animal-like promiscuity; moreover, it may be recognized wherever we consider not the individual but the overall dimension; not the person but the species; not the `form' but `matter' (in the Aristotelian sense of these two terms) … The will to equality is one and the same with the will to what is formless." (Men Among the Ruins) He argues extensively that an organic State with a unity of purpose will reject an imposed uniformity; that is, objects necessarily have separate forms, even though the ultimate purpose of these forms may be unitary.

I would interpret ``Consider[ing] an object in its all-sided wholeness and its inner bond with all other objects", on the other hand, to be the teaching that form is essentially borderless and that heterogeneity is an illusion. Most of the rest of this summary of Solovyov does not seem so conflicting, though. I will retract my comment about what is ``mechanical"; I forgot that in Evola, ``organic" is not synonymous with ``natural".


\hfill

\texttt{Alexander on 2012-11-30 at 20:10 said: }

``Consider[ing] an object in its all-sided wholeness and it's inner bond with all other objects" The archetypes pertain to the level of formless manifestation while principal unity is found beyond all manifestation.


\hfill

\texttt{Avery M on 2012-11-30 at 23:18 said: }

This seems to me to be the confusion that arises when language is insufficiently precise. In Guenon, ``all things are linked together by rigorous correspondences". Male corresponds to female, for instance. But male is not equivalent to female, which is the vibe I get from calling a correspondence a ``bond" and contrasting it with ``separateness".

To elaborate on my reference of Buddhism, the Heart Sutra's teaching of dependent arising and ``emptiness" can easily be misinterpreted to say that all things are homogenous, or that all things are equally valued, that is to say, valueless. Before we consider higher unity we must acknowledge separateness and difference in lower manifestation.


\hfill


\end{sffamily}\end{footnotesize}
